%------------------------------------------------------------------------------%
\documentclass[a4paper,12pt,fleqn]{article}

\usepackage{calculo_varias_variaveis-1}

\ShowAnswers

%------------------------------------------------------------------------------%
\begin{document}

\makeHeader{CFVVI}{Prova 3 -- Turma 6}{21/05/2026}

\bigskip
\textbf{Atenção}: a aplicação das regras do produto ou do quociente em casos
nos quais um dos termos é constante \textbf{será considerada incorreta},
por evidenciar compreensão inadequada das regras de derivação.

% 01
%------------------------------------------------------------------------------%
\exercise{50\%}
Considerando a função
\(
  f(x, y) = x^2y + e^{xy}
\)

\exerciseItem{20\%}
  Calcule o gradiente de $f$ em $(1, 0)$ e $(3, 4)$

\exerciseItem{30\%}
  Calcule a derivada direcional de $f$ no ponto
  $P=(1, 0)$, na direção do vetor $\vec{v}=(3, 4)$

\begin{answer}
\textbf{a)}
\vspace{-\baselineskip}
  \[
    f_x(x, y)
    = \D{}{x}\left(x^2y + e^{xy}\right)
    = 2xy + ye^{xy}
  \]
  \[
    f_y(x, y)
    = \D{}{y}\left(x^2y + e^{xy}\right)
    = x^2 + xe^{xy}
  \]
  \[
    \nabla f(x, y)
    = \vetor{f_x(x, y)}{f_y(x, y)}
    = \vetor{2xy + ye^{xy}}{x^2 + xe^{xy}}
  \]
  \[
    \nabla f(1, 0)
    = \vetor{2xy + ye^{xy}}{x^2 + xe^{xy}} \evalat{(1, 0)}{}
    = \vetor{2 \times 1 \times 0 + 0e^0}{1^2+1e^0}
    = \vetor{0}{2}
  \]
  \[
    \nabla f(3, 4)
    = \vetor{2xy + ye^{xy}}{x^2 + xe^{xy}} \evalat{(3, 4)}{}
    = \vetor{2 \times 3 \times 4 + 4e^{3\times 4}}
            {3^2 + 3e^{3 \times 4}}
    = \vetor{24+4e^{12}}{9+3e^{12}}
  \]

\bigskip
\textbf{b)}
\vspace{-\baselineskip}
  \[
    D_{\vec{u}}f(1, 0) = \nabla f(1, 0) \cdot \vec{u}
  \]
  O vetor $\vec{v}$ não é unitário. Sua norma é
  \(
    \norm{\vec{v}}
    = \sqrt{3^2 + 4^2}
    = 5
  \)

  Vetor unitário na direção de $\vec{v}$
  \[
    \vec{u}
    = \frac{\vec{v}}{\norm{\vec{v}}}
    = \Vetor{\frac{3}{5}}{\frac{4}{5}}
  \]
  Derivada direcionakl
  \[
    D_{\vec{u}}f(1, 0)
    = \nabla f(1, 0) \cdot \vec{u}
    = \vetor{0}{2} \cdot \Vetor{\frac{3}{5}}{\frac{4}{5}}
    = 0 \times \frac{3}{5} + 2  \times  \frac{4}{5}
    = \frac{8}{5}
  \]
\end{answer}

% 02
%------------------------------------------------------------------------------%
\clearpagequestiononly
\exercise{50\%}
Considerando a função
\(
  h(x, y) = \ln\left(x^2+y^2\right) + e^{xy}
\)
\exerciseItem{10\%} calcule seu gradiente
\exerciseItem{20\%} calcule a aproximação linear em torno de \((0, 2)\)
                    (use $\ln(4) = 1,39$)
\exerciseItem{20\%} use a aproximação linear para estimar \(f(0.05, 1.98)\)

\begin{answer}
\textbf{a)}
\vspace{-\baselineskip}
  \begin{align*}
    h_x(x, y)
    & = \D{}{x}\left(\ln\left(x^2+y^2\right) + e^{xy}\right) \\
    & = \frac{1}{x^2+y^2}\D{}{x}\left(x^2+y^2\right)
      + e^{xy}\D{}{x}\left(xy\right) \\
    & = \frac{2x}{x^2+y^2} + ye^{xy}
  \end{align*}
  \begin{align*}
    h_y(x, y)
    & = \D{}{y}\left(\ln\left(x^2+y^2\right) + e^{xy}\right) \\
    & = \frac{1}{x^2+y^2}\D{}{y}\left(x^2+y^2\right)
      + e^{xy}\D{}{y}\left(xy\right) \\
    & = \frac{2y}{x^2+y^2} + xe^{xy}
  \end{align*}
  \[
    \nabla h(x, y)
    = \vetor{h_x(x, y)}{h_y(x, y)}
    = \Vetor{\frac{2x}{x^2+y^2} + ye^{xy}}{\frac{2y}{x^2+y^2} + xe^{xy}}
  \]

\bigskip
\textbf{b)}
\vspace{-\baselineskip}
  \begin{align*}
    L(x, y)
    = h(a, b) + h_x(a, b)(x-a) + h_y(a, b)(y-b)
    = h(0, 2) + h_x(0, 2)x + h_y(0, 2)(y-2)
  \end{align*}
  \begin{align*}
    h(0, 2)
    = \left(\ln\left(x^2+y^2\right) + e^{xy}\right) \evalat{(0, 2)}{}
    = \ln\left(0+2^2\right) + e^{0}
    = \ln(4) + 1
    \approx 2,39
  \end{align*}
  \begin{align*}
    h_x(0, 2)
    = \left(\frac{2x}{x^2+y^2} + ye^{xy}\right) \evalat{(0, 2)}{}
    = \frac{0}{0+2^2} + 2e^0
    = 2
  \end{align*}
  \begin{align*}
    h_y(0, 2)
    = \left(\frac{2y}{x^2+y^2} + xe^{xy}\right) \evalat{(0, 2)}{}
    = \frac{2\times 4}{0+2^2} + 0e^{0}
    = \frac{8}{4}
    = 2
  \end{align*}
    \begin{align*}
      L(x, y)
      = h(0, 2) + h_x(0, 2)x + h_y(0, 2)(y-2)
      = 2,39 + 2x + 2(y-2)
    \end{align*}

\bigskip
\textbf{c)}
\vspace{-\baselineskip}
  \begin{align*}
    f(0.05, 1.98)
    \approx L(0.05, 1.98)
    & = \left(2,39 + 2x + 2(y-2)\right) \evalat{(0.05, 1.98)}{} \\
    & = 2,39 + 2\times 0,05 + 2(1,98-2) \\
    & = 2,39 + 0,10 - 0,04 \\
    & = 2,45
  \end{align*}
\end{answer}

%------------------------------------------------------------------------------%
\end{document}
%------------------------------------------------------------------------------%

